% dhsp.tex
%
% Talk for IQI Seminar, Caltech
%
% Pawel Wocjan <wocjan@cs.caltech.edu>

\documentclass[final,slideColor,nototal,colorBG,pdf,simple]{prosper}

%\input{Qcircuit.tex}

\hypersetup{pdfpagemode=FullScreen}

\usepackage[usenames]{color}
\usepackage{amsmath,amsfonts,amsthm} %,cancel}
\usepackage{pstricks,pst-node,pst-text,pst-3d,ulem}
\usepackage{epsfig}

\setlength{\abovedisplayskip}{-20pt}
\setlength{\belowdisplayskip}{-20pt}

%%%%%%%%%%%%%%%%%%%%%%%%%%%%%%%%%%%%%%%%%%%%%%%%%%%%%%%%%%%%%%%%%%%%%%%%%%%%%%%
% Macros

\definecolor{orange}{rgb}{1,.7,0}
\definecolor{gray}{rgb}{.5,.5,.5}
\definecolor{green}{rgb}{0,.7,0}

\newcommand{\bluetext}[1]{\textcolor{blue}{#1}}
\newcommand{\orangetext}[1]{\textcolor{orange}{#1}}
\newcommand{\greentext}[1]{\textcolor{green}{#1}}
\newcommand{\graytext}[1]{\textcolor{gray}{#1}}
\newcommand{\redbf}[1]{\textcolor{red}{\textbf{#1}}}
%\renewcommand{\CancelColor}{\red}

\newcommand{\bs}{\begin{slide}}
\newcommand{\es}{\end{slide}}

\newcommand{\be}{\begin{equation*}}
\newcommand{\ee}{\end{equation*}}
\def\ba#1\ea{\begin{align*}#1\end{align*}}
\newcommand{\nn}{\nonumber\\}

\newcommand{\<}{\langle} 
\renewcommand{\>}{\rangle}

\newcommand{\tr}{\operatorname{tr}}
\newcommand{\rank}{\operatorname{rank}}
\newcommand{\poly}{\operatorname{poly}}

\newcommand{\D}{{\mathcal D}}
\newcommand{\G}{{\mathcal G}}
\renewcommand{\H}{{\mathcal H}}

\newcommand{\C}{{\mathbb C}}
\renewcommand{\R}{{\mathbb R}}
\newcommand{\Z}{{\mathbb Z}}

\newcommand{\dg}{\dagger}

\newenvironment{mylist}{\begin{itemize}
                        \setlength{\itemsep}{2pt}}
                       {\end{itemize}}


 
%%%%%%%%%%%%%%%%%%%%%%%%%%%%%%%%%%%%%%%%%%%%%%%%%%%%%%%%%%%%%%%%%%%%%%%%%%%%%%%

\title{
\vskip -15pt
\textcolor{blue}{On the Quantum Hardness of}

\medskip 
\textcolor{blue}{Solving Isomorphism Problems as}

\medskip 
\textcolor{blue}{Nonabelian Hidden Shift Problems}}

\author{

\bigskip

\bigskip
\fontText{\textcolor{green}{Pawel Wocjan}}

\bigskip

\fontText{joint work with Andrew Childs}

\medskip 
\fontText{arXiv:quant-ph/0510185}}

%\institution{

%\bigskip
%\fontText{Computer Science Department \&

%Institute for Quantum Information

%\medskip
%California Institute of Technology}}

\slideCaption{}
%\Logo(9.5,7.5){\input{iqilogo2.eepic}}
\DefaultTransition{Replace}

\begin{document}
\maketitle

%%%%%%%%%%%%%%%%%%%%%%%%%%%%%%%%%%%%%%%%%%%%%%%%%%%%%%%%%%%%%%%%%%%%%%%%%%%%%%%

\bs{Organization of the talk}

Graph isomorphism

\bigskip
Hidden \textcolor{blue}{subgroup} approach

\bigskip
Hidden \textcolor{red}{shift} approach
\medskip
\begin{itemize}
\item Definition of nonabelian H-Shift-P

\medskip
\item Natural quantum approach to solving H-Shift-P

\medskip
\item Isomorphism problems as H-Shift-P

\bigskip
\item Structure of hidden shift states

\medskip
\item Single register measurement do not suffice 

\medskip
\item Lower bound on the number of copies
\end{itemize}


\es

%%%%%%%%%%%%%%%%%%%%%%%%%%%%%%%%%%%%%%%%%%%%%%%%%%%%%%%%%%%%%%%%%%%

\bs{Graph isomorphism}

$\Gamma$ graph on $n$ vertices and $g\in S_n$ permutation

\bigskip

$S_n$ acts on the set of $n$-vertex graphs:\quad $g(\Gamma)$ is the
graph obtained from $\Gamma$ by permuting the vertices according to
$g$

\bigskip
$\Gamma$ rigid $\Leftrightarrow$ $/\!\!\!\exists g\in S_n \, : \,
g(\Gamma)=\Gamma$

\bigskip
$\Gamma_0$ and $\Gamma_1$ isomorphic $\Leftrightarrow$ $\exists
s\in S_n\,:\, s(\Gamma_1)=\Gamma_0$

\medskip 
$s$ is the isomorphism between $\Gamma_1$ and $\Gamma_0$

\bigskip

\bigskip
Problem: given two rigid graphs decide if they are isomorphic or not
\es
%%%%%%%%%%%%%%%%%%%%%%%%%%%%%%%%%%%%%%%%%%%%%%%%%%%%%%%%%%%%%%%%%%%

\bs{Hidden subgroup problem}

$f : G \rightarrow S$ where $G$ is a group and $S$ a finite set

\medskip
$f$ hides a subgroup $H$, that is,
\begin{itemize}
\item $f$ is constant on the cosets $G/H$, and
\item $f$ is different on distinct cosets
\end{itemize}

\medskip
Problem: determine $H$
\es

%%%%%%%%%%%%%%%%%%%%%%%%%%%%%%%%%%%%%%%%%%%%%%%%%%%%%%%%%%%%%%%%%%%

\bs{Graph isomorphism as HSP}
\es

%%%%%%%%%%%%%%%%%%%%%%%%%%%%%%%%%%%%%%%%%%%%%%%%%%%%%%%%%%%%%%%%%%%

\bs{Quantum approach to solving HSP}

\es

%%%%%%%%%%%%%%%%%%%%%%%%%%%%%%%%%%%%%%%%%%%%%%%%%%%%%%%%%%%%%%%%%%%

\bs{Hardness of GI as HSP}

\es

%%%%%%%%%%%%%%%%%%%%%%%%%%%%%%%%%%%%%%%%%%%%%%%%%%%%%%%%%%%%%%%%%%%


\bs{Definition of nonabelian H-Shift-P}

black-box access to two injective functions 

\medskip
$f_0:G\rightarrow S$ and
$f_1:G\rightarrow S$, 

\medskip
where $G$ is a (nonabelian) group and $S$ is a finite set

\bigskip
the functions $f_0$ and $f_1$ are promised to satisfy:
\begin{itemize}
\item[(1)] there is an $s\in G$ such that $f_0(g)=f_1(gs)$ for all
$g\in G$

\medskip
$s$ is the hidden shift

\medskip
\item[(2)] the images of $f_0$ and $f_1$ are disjoint

\medskip
no hidden shift
\end{itemize}

\bigskip
Problem: determine if (1) or (2)
\es

%%%%%%%%%%%%%%%%%%%%%%%%%%%%%%%%%%%%%%%%%%%%%%%%%%%%%%%%%%%%%%%%%%%

\bs{Graph isomorphism as H-Shift-P}

$\Gamma$ graph on $n$ vertices and $g\in S_n$ permutation

\bigskip

\bigskip
$S_n$ acts on the set of $n$-vertex graphs:\quad $g(\Gamma)$ is the
graph obtained from $\Gamma$ by permuting the vertices according to
$g$

\bigskip

\bigskip
$\Gamma$ rigid $\Leftrightarrow$ $/\!\!\!\exists g\in S_n \, : \,
g(\Gamma)=\Gamma$

\bigskip

\bigskip
$\Gamma_0$ and $\Gamma_1$ isomorphic $\Leftrightarrow$ $\exists
s\in S_n\,:\, s(\Gamma_1)=\Gamma_0$

\bigskip 
$s$ is the isomorphism between $\Gamma_1$ and $\Gamma_0$
\es
%%%%%%%%%%%%%%%%%%%%%%%%%%%%%%%%%%%%%%%%%%%%%%%%%%%%%%%%%%%%%%%%%%%

\bs{Graph isomorphism as H-Shift-P}

\medskip
let $\Gamma_0$ and $\Gamma_1$ be two rigid $n$-vertex graphs 

\bigskip
define $f_0(g)=g(\Gamma_0)$ and $f_1(g)=g(\Gamma_1)$

\bigskip

\bigskip
clearly, $f_0$ and $f_1$ are injective because both graphs rigid

\bigskip
\begin{itemize}
\item[(1)] assume $s(\Gamma_1)=\Gamma_0$ for $s\in S_n$ \quad
$\Rightarrow$ \quad $f_0(g)=f_1(gs)$

\bigskip
\item[(2)] assume $\Gamma_1$ and $\Gamma_0$ are not
isomorphic \quad $\Rightarrow$ 

\medskip
the images of $f_0$ and $f_1$ are disjoint
\end{itemize}

\es

%%%%%%%%%%%%%%%%%%%%%%%%%%%%%%%%%%%%%%%%%%%%%%%%%%%%%%%%%%%%%%%%%%%
\bs{Quantum approach to solving H-Shift-P}

prepare a uniform superposition over $i\in\mathbb{Z}_2$ and $g\in G$,
compute the value of $f_i(g)$, and measure the third register
\[
\frac{1}{\sqrt{2|G|}}\sum_{g\in G} 
(|0,g,f_0(g)\> + |1,g,f_1(g)\>)
\]

\begin{itemize}
\item if there is a hidden shift {\red $s$}, then we are left with
\[
|\phi_{{\red s},g}\>:=\frac{1}{\sqrt{2}}(|0,g\> + |1,g{\red s}\>)
\quad\quad\in\C^2\otimes\C^{|G|}
\]
for some uniformly random (unknown) $g\in G$

\medskip
\item if there is no hidden shift, then we obtain
\[
|i,g\>\quad\quad\in\C^2\otimes\C^{|G|}
\]
for uniformly random $i\in\mathbb{Z}_2$ and $g\in G$
\end{itemize}
\es

%%%%%%%%%%%%%%%%%%%%%%%%%%%%%%%%%%%%%%%%%%%%%%%%%%%%%%%%%%%%%%%%%%%
\bs{Hidden shift $s$}
if $s$ is the hidden shift, then the density matrix is
\[
\gamma_1(s):=\frac{1}{|G|}\sum_{g\in G} |\phi_{s,g}\>\<\phi_{s,g}|
\]

\bigskip
Observation:\quad $\gamma_1(s)$ can be expressed as 
$R$ 
\[
\gamma_1(s)
= 
\frac{1}{2|G|}
\left(
\begin{array}{cc}
I         & R(s) \\
R(s^{-1}) & I
\end{array}
\right)
\]

\bigskip
where $R$ is the right regular representation 
\[
R(s)|g\> = |gs^{-1}\>
\]
\es


%%%%%%%%%%%%%%%%%%%%%%%%%%%%%%%%%%%%%%%%%%%%%%%%%%%%%%%%%%%%%%%%%%%
\bs{No hidden shift}
if there is no hidden shift, the density matrix is
\[
\gamma_2:=
\frac{1}{2|G|}\sum_{i\in\mathbb{Z}_2}\sum_{g\in G} 
|i,g\>\<i,g|
\]

\bigskip
$\gamma_2$ is the maximally mixed state
\[
\gamma_2 := 
\frac{1}{2|G|}
\left(
\begin{array}{cc}
I & 0 \\
0 & I
\end{array}
\right)
\]
\es


%%%%%%%%%%%%%%%%%%%%%%%%%%%%%%%%%%%%%%%%%%%%%%%%%%%%%%%%%%%
\bs{H-Shift-P as state discrimination}
we can apply the above procedure $k$ times to obtain $k$ copies of the
hidden shift state (or the maximally mixed state)

\bigskip
assume that in the case where there is a hidden shift $s$, it is
equally likely to correspond to any element of $G$

\bigskip
then the problem is to distinguish the two density operators
\begin{eqnarray*}
\gamma_1^{(k)} & := & \frac{1}{|G|} \sum_{s\in G} \gamma_1^{(k)}(s) \\
\gamma_2^{(k)} & := & \frac{1}{(2|G|)^2} I
\end{eqnarray*}

\bigskip
where $\gamma_1^{(k)}(s):=\gamma_1(s)^{\otimes k}$
\es

%%%%%%%%%%%%%%%%%%%%%%%%%%%%%%%%%%%%%%%%%%%%%%%%%%%%%%%%%%%%%%%%%%%%%%%%%%%%%%

\bs{Structure of the hidden shift states}
Recall
\[
\gamma_1(s) = \frac{1}{2|G|} \left(
\begin{array}{cc}
I         & R(s) \\
R(s^{-1}) & I
\end{array}
\right)
\]
where $R$ the right regular representation $R(s)|g\>=|gs^{-1}\>$

\bigskip
$\hat{G}$ a set of all irreducible (inequivalent)
representations of $G$

\medskip
$F$ the (corresponding) Fourier transform over $G$

\bigskip
then we have block decomposition
\[
F R(s) F^\dagger = \bigoplus_{\rho\in\hat{G}} I_{d_\rho} \otimes \rho(s)
\]

\es

%%%%%%%%%%%%%%%%%%%%%%%%%%%%%%%%%%%%%%%%%%%%%%%%%%%%%%%%%%%%%%%%%%%%%%%%%%%%%%

\bs{Structure of the hidden shift states}
in the Fourier basis we have the block decomposition

\bigskip
\begin{eqnarray*}
\tilde{\gamma}_1(s) & = & \frac{1}{2|G|} 
\bigoplus_{\rho\in\hat{G}}
I_{d_\rho} \otimes
\underbrace{
\left(
\begin{array}{cc}
I_{d_\rho}   & \rho(s) \\
\rho(s^{-1}) & I_{d_\rho}
\end{array}
\right)}_{\mbox{{\large $B^\rho(s)$}}} \\
%%%%%%%%%%%%%%%%%%%%%%%%%%%%%%%%%%%%%%%%%%%%%%%%%%%%%%%%%%%%%%%%%%%
\tilde{\gamma}_2 & := & 
\frac{1}{2|G|} 
\bigoplus_{\rho\in\hat{G}}
I_{d_\rho} \otimes
\underbrace{\left(
\begin{array}{cc}
I_{d_\rho} & 0 \\
0          & I_{d_\rho}
\end{array}
\right)}_{\mbox{{\large $I_{2d_\rho}$}}}
\end{eqnarray*}

\es

%%%%%%%%%%%%%%%%%%%%%%%%%%%%%%%%%%%%%%%%%%%%%%%%%%%%%%%%%

\bs{Single-register measurements}

without loss of generality any POVM used to distinguish among these
states has the same block structure
\[
\mbox{{\LARGE $\Downarrow$}}
\]
\medskip
first measure representation name (weak Fourier sampling), then
perform arbitrary POVM inside the block

\bigskip

\bigskip
the various irreducible representations $\rho$ occur independently
according to the Plancherel distribution
\[
{\rm Pr}(\rho) = \frac{d_\rho^2}{|G|}
\]
regardless of whether or not there is a hidden shift

\es

%%%%%%%%%%%%%%%%%%%%%%%%%%%%%%%%%%%%%%%%%%%%%%%%%%%%%%%%%%%%%%%%%%%%%%%%
\bs{Single register measurements} 

${\cal E}=\{a_1|\psi_1\>\<\psi_1|,\ldots,a_r|\psi_r\>\<\psi_r|\}$ 
arbitrary POVM acting on subspace of dimension $2d_\rho$ corresponding
to observed $\rho$

\bigskip
no hidden shift $\rightarrow$ 

post-measurement state is $I_{2d_\rho}/2d_\rho$ and probability of $j$ is
\[
p_2(j)=\frac{a_j}{2d_\rho} \<\psi_j|I_{2\rho_j}|\psi_j\>
= \frac{a_j}{2d_\rho}
\]

\bigskip
hidden shift $s$ $\rightarrow$ 

post-measurement state is $B^\rho(s)/2d_\rho$ and probability of $j$ is
\[
p_1(j|s) = \frac{a_j}{2d_\rho} \<\psi_j|B^\rho(s)|\psi_j\>
\]

\bigskip
let $P_2$ and $P_{1,s}$ corresponding probability distributions

\es

%%%%%%%%%%%%%%%%%%%%%%%%%%%%%%%%%%%%%%%%%%%%%%%%%%%%%%%%%%%

\bs{Single register measurements}
single register measurement do not suffice because

\[
(*)\quad
{\rm Pr}_{s\in G,\rho\in\hat{G}} (\|P_{1,s}-P_2\|_{\rm TV} \ge e^{-\Theta(n)})
\le e^{-\Theta(n)}
\]

\bigskip
define $p_1(j):=\frac{1}{|G|}\sum_{s\in G} p_1(j|s)$ 

let $P_1$ be the corresponding probability distribution

\bigskip
we prove $(*)$ by showing
\begin{itemize}
\item $P_1$ and $P_2$ are equal for all
$\rho\in\hat{G}-\{\hat{1}\}$
\item ${\rm Pr}_{s\in G,\rho\in\hat{G}} (\|P_{1,s}-P_1\|_{\rm TV} 
\ge e^{-\Theta(n)}) \le e^{-\Theta(n)}$
\end{itemize}
\es

%%%%%%%%%%%%%%%%%%%%%%%%%%%%%%%%%%%%%%%%%%%%%%%%%%%%%%%%%%%%%%%%%%

\bs{Single register measurements}
let $\rho\in\hat{G}$ be fixed and $d_\rho$ be its degree

\bigskip
let $\sigma_j$ be the variance of the random variable $p_1(j|s)$ with
respect to $s\in G$ uniformly at random
\[
\sigma_j^2:=\frac{1}{|G|}\sum_{s\in G} p_1(j|s)^2-p_1(j)^2
\]

Chebyshev's inequality
\[
{\rm Pr}_{s\in G}(|p_1(j|s)-p_1(j)|
\ge a_j c)\le\frac{\sigma_j^2}{a_j^2 c^2}
\]

define 
\[
J^s_{\rm bad}:=\{j\in J\,:\,|p_1(j|s)-p_1(j)|\ge a_j c\} \mbox{ and }
J^s_{\rm good}:=J-J^s_{\rm bad}
\]
\es

%%%%%%%%%%%%%%%%%%%%%%%%%%%%%%%%%%%%%%%%%%%%%%%%%%%%%%%%%%%%%%%%%%%%%%
\bs{Single register measurements}
Now we can write 
\[
\|P_{1,s}-P_1\|_{\rm TV}=
\sum_{j\in J^s_{\rm good}} |p_1(j|s)-p_1(j)| +
\sum_{j\in J^s_{\rm bad}}  |p_1(j|s)-p_1(j)|
\]

\bigskip
Markov's inequality gives the bound
\[
{\rm Pr}(\sum_{j\in J^s_{\rm bad}} a_j\ge c') \le 
\frac{1}{d_\rho^2 c^2 c'}
\]
\es

%%%%%%%%%%%%%%%%%%%%%%%%%%%%%%%%%%%%%%%%%%%%%%%%%%%%%%%%%%%%%%%%%%%

\bs{Lower number of copies}

\es

%%%%%%%%%%%%%%%%%%%%%%%%%%%%%%%%%%%%%%%%%%%%%%%%%%%%%%%%%%%%%%%%%%%

\end{document}

